% Proved after full independent audits. Potential CM means finite-extension endomorphisms.
\subsection{The last Lubin--Tate layer and the first wild CM index}

Use the definition of potentially CM in B52: the rational
endomorphism algebra over some finite extension with good
reduction contains a quadratic field over \(\mathbf Q_p\).
This is not the endomorphism algebra of an abstract constant
generic etale group over an algebraic closure.

\begin{lemma}[The nonsquare integral CM model at arbitrary index]
\label{lem:B-RD78-model}
\status{Proved; complete independent root and companion proof, primary-source and exact-code audits passed.}
Let \(p\ge5\), and let \(K\) be complete discretely valued
of mixed characteristic \((0,p)\), with algebraically closed
residue field and \(\mu_p\subset K\). There is no restriction
on its absolute index \(e\). Let \(M\) be a potentially CM,
potentially supersingular height-two group, whose determinant
is \(\chi_p\psi\), with nonsquare finite
\(\mu_{p-1}\)-valued character \(\psi\).
Then its quadratic endomorphism field is the unramified field
\(E=\mathbf Q_{p^2}\); all finite-extension endomorphisms are
defined over \(K\), and its Tate lattice is a free rank-one
\(\mathcal O_E\)-module. After choosing its tangent orientation,
\[
        T_p(M)=\mathcal O_E(\theta\epsilon),          \tag{B-RD78.1}
\]
where \(\theta\) is the character of the standard formal
\(\mathcal O_E\)-module \([p](T)=pT+T^{p^2}\), base changed
to \(K\), and \(\epsilon\) is finite and tame with values
in \(\mu_{p^2-1}\).
\end{lemma}
\begin{proof}
We verify the dependence on the absolute index in the complete
proofs of Lemmas~\ref{lem:B-RD52-CM} and~\ref{lem:B-RD52-LT}.
The faithful good extension is tame of degree dividing
\(p^2-1\) by Lemma~\ref{lem:B-RD49-inertia}, at arbitrary
\(e\). Tate full faithfulness extends integral endomorphisms
to its good model; clearing a power of \(p\) does the same
for rational ones. Reduction is injective by
\cite[Proposition~4.2 and proof]{HopkinsGross1994Bundles}.
Since the residue field is algebraically closed, finite-extension
endomorphisms therefore descend to the good field. Their
rational algebra is a division algebra acting on a two-dimensional
\(\mathbf Q_p\)-space, and the CM assumption makes it a
quadratic field.

If that field is ramified, its special prime-to-\(p\) normalizer
has order dividing \(2(p-1)\). For \(p\equiv3\pmod4\)
its reduced norms are squares. For \(p\equiv1\pmod4\), the
generic normalizer consists of multiplication by a unit and
such multiplication followed by conjugation; its determinants
are \(\pm N_{E/\mathbf Q_p}(u)\), both squares modulo \(p\).
Here \(\chi_p\equiv1\pmod p\), so either case contradicts
the nonsquare \(\psi\). Thus \(E\) is unramified.
The noncentral coset of its normalizer in the special division
algebra consists of \(E^\times\Pi\), whose valuations are
odd, so has no units. Special finite inertia centralizes \(E\),
and reduction injectivity makes the generic action centralize
\(E\) as well. This also descends the endomorphisms to \(K\).

Choose a Galois element with nonsquare finite determinant.
Its action is multiplication by a unit of \(E\) with nonsquare
residue norm, so its residue does not lie in \(\mathbf F_p\).
The algebra generated by that action over \(\mathbf Z_p\)
is \(\mathcal O_E\), since its quadratic polynomial has
irreducible separable reduction. The stable Tate lattice is
therefore free of rank one over \(\mathcal O_E\).

The integral action extends to the good model by Tate full
faithfulness. The explicit unit conjugator and deformation
argument in Lemma~\ref{lem:B-RD52-LT} identify this model with
the standard oriented Lubin--Tate group: averaging the universal
parameter against a primitive \((p^2-1)\)st-root action gives
a parameter multiplied by \(\lambda^{p-1}\ne1\), and an
\(\mathcal O_E\)-linear deformation must send it to zero.
The quotient of its generic character by \(\theta\) is trivial
on the tame good field, hence is the finite tame character
\(\epsilon\) in (B-RD78.1). None of these steps uses
\(p\nmid e\). Only the subsequent first-wild-layer claim
of B52 used that hypothesis; we replace it below.
\end{proof}

\begin{lemma}[A full terminal layer after wild base change]
\label{lem:B-RD78-layer}
\status{Proved; complete independent root and companion proof, primary-source and exact-code audits passed.}
Retain the preceding notation, set \(v=v_p(e)\), and let
\(L/K\) be any common good tame extension. If \(n\ge v+2\),
the image of \(\theta|_{G_L}\) modulo \(p^n\) contains
\[
          (1+p^{n-1}\mathcal O_E)/(1+p^n\mathcal O_E).
                                                        \tag{B-RD78.2}
\]
\end{lemma}
\begin{proof}
Choose a coefficient ring \(W(k)\subset\mathcal O_K\)
by the Cohen structure theorem, and set \(K_0=W(k)[1/p]\).
This uses \cite[Theorem~10.160.8, Tag~032A]{Stacks032A}.
Because the residue field is perfect and \(K\) is a complete
discrete valuation field, its ring of integers is finite free
of rank \(e\) over this coefficient ring: a uniformizer and
Weierstrass preparation give an Eisenstein relation of degree
\(e\), with basis \(1,\pi,\ldots,\pi^{e-1}\).
Thus \([L:K_0]=e[L:K]\), whose \(p\)-adic valuation is \(v\).
No identification of \(k\) with \(\overline{\mathbf F}_p\)
is made.

Let \(D_n\) be the full \(p^n\)-division field of the standard
Lubin--Tate group over \(K_0\), and put \(q=p^2\).
The direct proof of
\cite[Theorem~2(b),(d),(e) and complete proof, pp.~383--384]
{LubinTate1965FormalCM} works after this residue extension.
Explicitly, the primitive division polynomial
\[
 \Phi_n(T)=\bigl([p]^{\circ(n-1)}(T)\bigr)^{q-1}+p
                                                        \tag{B-RD78.3}
\]
is Eisenstein of degree \((q-1)q^{n-1}\), and a primitive
point generates the free rank-one \(\mathcal O_E/p^n\)-module
of torsion. Every other point is its \(\mathcal O_E\)-multiple.
The faithful Galois injection into
\((\mathcal O_E/p^n)^\times\), together with that Eisenstein
degree, is therefore an equality, and \(D_n/K_0\) is totally
ramified. This is an elementary division-field argument, not
an appeal to local class field theory over the larger residue
field.

Set \(U=(1+p\mathcal O_E)/(1+p^n\mathcal O_E)\), and let
\(J\) be the image after base change to \(L\). Galois
intersection gives
\[
 [\operatorname{Gal}(D_n/K_0):J]=[D_n\cap L:K_0],
 \qquad [U:U\cap J]\le p^v.                         \tag{B-RD78.4}
\]
The second index is the \(p\)-part of the first, since the
unit group is abelian with prime-to-\(p\) Teichmuller factor;
the intersection degree divides \([L:K_0]\).
An abelian group of order at most \(p^v\) is killed by \(p^v\),
so \(U^{p^v}\subset J\). For odd \(p\), successive binomial
expansions give
\[
 U^{p^v}=(1+p^{v+1}\mathcal O_E)/(1+p^n\mathcal O_E).
                                                        \tag{B-RD78.5}
\]
Indeed the \(p\)th-power map has leading term
\((1+p^r b)^p=1+p^{r+1}b\pmod{p^{r+2}}\) for \(r\ge1\);
successive lifts give surjectivity onto the next layer.
For \(n\ge v+2\), (B-RD78.5) contains (B-RD78.2).
\end{proof}

\begin{proposition}[Nonsquare CM character recovery]
\label{prop:B-RD78-CM}
\status{Proved; complete independent root and companion proof, primary-source and exact-code audits passed.}
Let \(K\) have perfect residue field and the other hypotheses
above, and put \(v=v_p(e)\). Two potentially CM, potentially
supersingular height-two groups with common nonsquare finite
determinant and isomorphic actual \(p^n\)-torsion have equal
special inertia characters whenever \(n\ge v+2\).
\end{proposition}
\begin{proof}
Pass to the completed maximal unramified field and choose
compatible orientations. The first lemma gives
\(T_p(M_i)=\mathcal O_E(\theta\epsilon_i)\), with the same
\(\theta\). On a common good tame field, both finite twists
are trivial. An intertwiner modulo \(p^n\) commutes with every
operator \(1+p^{n-1}b\), for \(b\in\mathcal O_E\), by the
preceding lemma. Dividing the commutator by \(p^{n-1}\) shows
that its reduction modulo \(p\) centralizes all multiplication
operators from \(\mathbf F_{p^2}\). This centralizer consists
of field multiplication itself. The reduction of the intertwiner
is therefore multiplication by a nonzero field element.

The full intertwining identity for every \(g\in G_K\), reduced
modulo \(p\), now gives
\(\bar\epsilon_1(g)=\bar\epsilon_2(g)\), after cancelling
that scalar and the nonzero scalar \(\bar\theta(g)\).
Reduction is injective on \(\mu_{p^2-1}\), so the two finite
characters agree. The good factor has trivial finite special
inertia; the covariant dictionary of B52 therefore gives the
same requested character for the two groups.
\end{proof}

\begin{theorem}[An exact first-wild index in the potential-CM subclass]
\label{thm:B-RD78}
\status{Proved; complete independent root and companion proof, primary-source and exact-code audits passed.}
For \(p\ge5\), \(e=p(p-1)\), \(\mu_p\subset K\), and
perfect residue field, actual \(p^3\)-torsion determines the
special character of two potentially CM, potentially supersingular
height-two groups. Actual \(p^2\)-torsion does not always do so
within this same subclass. Its least manuscript index is thus
exactly two.

Consequently any remaining full B.4 counterexample at actual
\(125\)-torsion for \(p=5,e=20\) must have at least one
whole formal model without quadratic finite-extension
endomorphisms, in either normalized pattern of B77.
\end{theorem}
\begin{proof}
For the upper bound, \(v_p(e)=1\), and the finite determinants
agree by Lemma~\ref{lem:B-RD76-determinant}. A square finite
determinant is covered by Proposition~\ref{prop:B-RD77-square};
the nonsquare case is the preceding proposition with \(n=3\).

For a lower example take \(k=\overline{\mathbf F}_p\),
\(K_0=W(k)[1/p]\), and the standard Lubin--Tate group \(F_0\).
Let \(D_2/K_0\) be its full \(p^2\)-division field, with
\[
 \operatorname{Gal}(D_2/K_0)
 = (\mathcal O_E/p^2)^\times
 =\mu_N\times(1+p\mathcal O_E)/(1+p^2\mathcal O_E),
 \quad N=p^2-1.
                                                        \tag{B-RD78.6}
\]
In this actual Galois group take
\[
 H=\mu_{p+1}\times
       \{1+pa:a\in\mathbf F_p\},\qquad K=D_2^H.
                                                        \tag{B-RD78.7}
\]
The field is completely specified. Its absolute ramification is
\(Np^2/[p(p+1)]=p(p-1)\). The determinant of the good
height-two dimension-one group \(F_0\) is cyclotomic and equals
the norm of \(\theta\). The subgroup \(H\) has norm one
modulo \(p\), hence \(\mu_p\subset K\).

Choose a uniformizer \(\pi\) of \(K\), and let
\(\omega:G_K\to\mu_N\) be the tame character from
\(\pi^{1/N}\). The first division equation is \(T^N=-p\).
Since the unit in \(-p=\pi^e u\) has an \(N\)th root in
\(K\) by Hensel's lemma, its Teichmuller torsion character is
\(\bar\theta=\bar\omega^e\). Write \(\sigma\) for the
unramified automorphism of \(\mathcal O_E\).
On the exact mod-\(p^2\) image \(H\), the wild factor in
(B-RD78.7) is scalar and fixed by \(\sigma\). Therefore
\[
 \frac{\sigma(\theta)}{\theta}
       =\omega^{(p-1)e}\pmod{p^2}.
                                                        \tag{B-RD78.8}
\]
Define actual finite CM twists of \(F_0\) by
\[
 \epsilon_1=\omega,\qquad
 \epsilon_2=\omega^{3p-2}.
                                                        \tag{B-RD78.9}
\]
They are effective generic BT twists, and acquire the original
good CM model over a tame extension. Since
\(p+(p-1)e\equiv3p-2\pmod N\), equation (B-RD78.8) says
that \(\sigma\), as a \(\mathbf Z/p^2\)-linear map of
\(\mathcal O_E/p^2\), intertwines
\(\theta\epsilon_1\) with \(\theta\epsilon_2\).
This is an isomorphism of the actual \(p^2\)-torsion modules,
and hence of generic finite group schemes.

Their finite covariant special eigenvalue exponents are
\[
                 \{1,p\},\qquad \{3p-2,3-2p\}
                 \pmod N.                            \tag{B-RD78.10}
\]
Here the good factor has trivial finite descent and the finite
twist contributes \(\epsilon_i\), exactly as in the explicit
covariant convention of B52. The exponent \(3p-2\) is neither
\(1\) nor \(p\) modulo \(N\): their differences are
\(3(p-1)\) and \(2(p-1)\), and \(p+1\ge6\) divides
neither \(3\) nor \(2\). The characters therefore differ.
Both finite norms are the same primitive character, since
\(3p-2\equiv1\pmod{p-1}\). This supplies the lower example
inside the stated CM subclass. At \(p=5\) its finite types
are \(A\) and \(C\).
\end{proof}

\begin{remark}[The unrestricted wild problem is not replaced by CM]
\label{rem:B-RD78-scope}
\status{Proved; complete independent root and companion proof, primary-source and exact-code audits passed.}
The terminal-layer proposition uses an index bound on an actual
Lubin--Tate division image and integral endomorphisms; it does
not assert that an arbitrary height-two formal group is CM.
The lower pair is an actual finite extension and actual pair
of generic BT twists, not merely two compatible finite characters.
The first-wild full least index remains in \(\{2,3\}\) until
the remaining non-CM \(p^3\)-comparisons are resolved.


\smallskip
\noindent\emph{Retained exact supplement.}
The independently rerunnable program
\path{code/reduction_descent/verify_wild_CM_layer.py}
checks 15,696 actual mod-\(p^2\) twist identities in explicit
unramified quadratic rings, 57,906 bounded-index additive
subgroups and their 115,812 nonzero terminal-layer basis tests,
46,228 residual centralizer matrices, 24 primitive Eisenstein
reductions, and 24 prime-index character rows.
The corrected author run passed in
\(0.21962125301070046\) seconds.
The first executed version and output are retained separately:
its additional terminal test used the ambient zero class
\(p^{n-1}\), although its stronger \(p^v\)-containment test
was already correct. The corrected explicit coordinate is
\(p^{n-2}\), because logarithmic coordinates for
\((1+p\mathcal O_E)/(1+p^n\mathcal O_E)\) divide by \(p\).
No finite range replaces the field construction or the
all-prime proof above.
The independent root and companion audits read every proof and
the entire corrected program; unchanged-program reruns passed in
\(0.21194225999352057\) and \(0.22546726299333386\) seconds.
Both freshly read the original Lubin--Tate proofs on
pp.~381--384 and the complete Stacks Tag~032A proof, retaining
the finite-residue limitation of the original Lubin--Tate source.
The exact audited candidate and every executed earlier output
remain preserved.
\end{remark}
